%----------------------------------------------------------------------------------------
% SCALAIR PRODUCT, HOEKEN EN AFSTANDEN. Alle opgaven zijn eigen werk.
%
% Opbouw: Start, leerstof (theorie, werkwijze of uitleg, voorbeeld, korte oefening), en op een
% nieuwe pagina de Opdrachten (door elkaar, stijgende moeilijkheid), Koppel en Test jezelf.
% Labels van dit hoofdstuk beginnen met h4-.
%----------------------------------------------------------------------------------------

\mychapter{4}{Hoeken en afstanden}

\lkr{Leerplan B+S'': LPD 5-8 (ruimtemeetkunde), het deel over het scalair product, loodrechte stand, hoeken en afstanden. Het scalair product in het vlak kennen sommige leerlingen uit de 2de graad; hier wordt het in de ruimte ingevoerd en meteen gebruikt voor de normaalvector van een vlak. Dat is de sleutel tot de rest van het hoofdstuk.}

Hoe bereken je de hoek tussen twee ribben van een piramide, of de afstand van een punt tot een vlak? Met coördinaten gaat dat zonder figuur, dankzij één bewerking: het scalair product van twee vectoren. Het zegt meteen of twee vectoren loodrecht op elkaar staan.

\begin{oefeningen}[Start]
\oef{h4-start-1}

\oef{h4-verken}
\end{oefeningen}

%========================================================================================
\section{Het scalair product}

\begin{theorie}
Het \textbf{\textcolor{accent}{scalair product}} van $\vec v = \vcr{v_1}{v_2}{v_3}$ en $\vec w = \vcr{w_1}{w_2}{w_3}$ is het getal
\[\formule{\vec v \cdot \vec w = v_1w_1 + v_2w_2 + v_3w_3}.\]
\begin{stelling}{Scalair product en hoek}
Is $\theta$ de hoek tussen $\vec v$ en $\vec w$ (beide niet $\vec 0$), dan is
\[\formule{\vec v \cdot \vec w = \norm{\vec v}\,\norm{\vec w}\cos\theta}.\]
In het bijzonder: $\formule{\vec v \perp \vec w \iff \vec v \cdot \vec w = 0}$.
\end{stelling}
\begin{bewijs}
Neem de driehoek met zijden $\vec v$, $\vec w$ en $\vec v - \vec w$. De cosinusregel geeft
\bewijsregel{\norm{\vec v - \vec w}^2 = \norm{\vec v}^2 + \norm{\vec w}^2 - 2\norm{\vec v}\norm{\vec w}\cos\theta.}
Met coördinaten is $\norm{\vec v - \vec w}^2 = \sum (v_i - w_i)^2 = \norm{\vec v}^2 + \norm{\vec w}^2 - 2(v_1w_1 + v_2w_2 + v_3w_3)$. Vergelijk beide.
\end{bewijs}
\end{theorie}

\voorbeeld{}{Bereken de hoek tussen $\vec v = \vcr{1}{2}{2}$ en $\vec w = \vcr{2}{-1}{2}$.\label{vb:h4-hoek}}{%
\uitleg{Bereken het scalair product en de normen, en los op naar $\cos\theta$.}
\begin{stappenbreed}
\vec v \cdot \vec w = 2 - 2 + 4 = 4,\quad \norm{\vec v} = \norm{\vec w} = 3 & \\
\cos\theta = \tfrac{4}{9} \iff \theta \approx 63{,}6^\circ & Rekentoestel in graden.
\end{stappenbreed}}

\begin{oefeningen}
\oef{h4-hoek}
\end{oefeningen}

\voorbeeld{}{Welke hoek maakt een lichaamsdiagonaal van een kubus met een ribbe?\label{vb:h4-kubus}}{%
\uitleg{Neem de kubus met ribbe 1 in de oorsprong. De diagonaal is $\vcr{1}{1}{1}$, de ribbe $\vcr{1}{0}{0}$.}
\begin{stappenbreed}
\cos\theta = \frac{1}{\sqrt3 \cdot 1} \iff \theta \approx 54{,}7^\circ & Voor elke kubus dezelfde hoek.
\end{stappenbreed}}

\begin{oefeningen}
\oef{h4-scalair}

\oef{h4-kubushoek}
\end{oefeningen}

%========================================================================================
\section{Normaalvector en loodrechte stand}

\naast{%
\begin{theorie}
Een \textbf{\textcolor{accent}{normaalvector}} van een vlak staat loodrecht op elke vector in dat vlak. Het vlak $\formule{ax + by + cz = d}$ heeft normaalvector $\formule{\vec n = \vcr{a}{b}{c}}$.

Het vlak door $P_0(x_0, y_0, z_0)$ met normaalvector $\vcr{a}{b}{c}$ is $\formule{a(x - x_0) + b(y - y_0) + c(z - z_0) = 0}$.
\end{theorie}}{%
\begin{tikzpicture}[ruimte,scale=.7,font=\scriptsize]
  \fill[accent!12] (0,0,0) -- (3,0,0) -- (3,3,0) -- (0,3,0) -- cycle;
  \draw[accent] (0,0,0) -- (3,0,0) -- (3,3,0) -- (0,3,0) -- cycle;
  \draw[->,accent,very thick] (1.5,1.5,0) -- (1.5,1.5,2) node[above,text=black] {$\vec n$};
  \draw[->,grijs,thick] (1.5,1.5,0) -- (1.5,2.8,0);
  \draw (1.5,1.5,.3) -- (1.5,1.8,.3) -- (1.5,1.8,0);
\end{tikzpicture}}

Waarom? Liggen $P_0$ en $P$ in het vlak, dan is $a(x - x_0) + b(y - y_0) + c(z - z_0) = d - d = 0$: dat is $\vec n \cdot \overrightarrow{P_0P} = 0$. Met de normaalvector beslis je over loodrechte stand:
\begin{itemize}
\item een rechte staat loodrecht op een vlak als haar richtingsvector collineair is met $\vec n$;
\item twee vlakken staan loodrecht op elkaar als $\vec n_1 \cdot \vec n_2 = 0$;
\item de hoek tussen twee vlakken is de (scherpe) hoek tussen hun normaalvectoren.
\end{itemize}

\voorbeeld{}{Geef een vergelijking van het vlak door $A(1, 2, 3)$ dat loodrecht staat op de rechte $l\colon (0, 0, 1) + r\vcr{2}{-1}{1}$.\label{vb:h4-loodvlak}}{%
\uitleg{De richtingsvector van $l$ is een normaalvector van het gezochte vlak.}
\begin{stappenbreed}
2(x - 1) - (y - 2) + (z - 3) = 0 & \\
2x - y + z = 3 & Controle: $A$ geeft $2 - 2 + 3 = 3$.
\end{stappenbreed}}

\begin{oefeningen}
\oef{h4-loodvlak}
\end{oefeningen}

\voorbeeld{}{Staan de vlakken $x + 2y - z = 4$ en $3x - y + z = 0$ loodrecht op elkaar?\label{vb:h4-loodvlakken}}{%
\uitleg{Bereken het scalair product van de normaalvectoren.}
\begin{stappenbreed}
\vcr{1}{2}{-1} \cdot \vcr{3}{-1}{1} = 3 - 2 - 1 = 0 & Ja: loodrecht.
\end{stappenbreed}}

\begin{oefeningen}
\oef{h4-normaal}

\oef{h4-hoekvlakken}
\end{oefeningen}

%========================================================================================
\section{Afstanden}

\begin{theorie}
\begin{stelling}{Afstand van een punt tot een vlak}
De afstand van $P(x_0, y_0, z_0)$ tot het vlak $ax + by + cz = d$ is
\[\formule{d(P, \alpha) = \frac{|ax_0 + by_0 + cz_0 - d|}{\sqrt{a^2 + b^2 + c^2}}}.\]
\end{stelling}
\end{theorie}

De afstand is de lengte van de loodlijn uit $P$ op het vlak. Die loodlijn heeft richting $\vec n$; de formule volgt door haar snijpunt met het vlak te berekenen (opdracht~\ref{oef:h4-o8}). De afstand van een punt tot een rechte vind je met een loodvlak: het vlak door $P$ loodrecht op de rechte snijdt ze in het voetpunt van de loodlijn.

\voorbeeld{}{Bereken de afstand van $P(1, 2, 3)$ tot het vlak $2x - y + 2z = 1$.\label{vb:h4-afstandvlak}}{%
\uitleg{Vul $P$ in het linkerlid in, trek $d$ af, en deel door de norm van $\vec n$.}
\begin{stappenbreed}
\frac{|2 - 2 + 6 - 1|}{\sqrt{4 + 1 + 4}} = \frac{5}{3} &
\end{stappenbreed}}

\begin{oefeningen}
\oef{h4-afstandvlak}
\end{oefeningen}

\voorbeeld{}{Bereken de afstand van $P(3, 1, 2)$ tot de rechte $l\colon r\vcr{1}{1}{1}$ door de oorsprong.\label{vb:h4-afstandrechte}}{%
\uitleg{Het vlak door $P$ loodrecht op $l$ heeft normaalvector $\vcr{1}{1}{1}$. Zijn snijpunt met $l$ is het voetpunt $V$.}
\begin{stappenbreed}
x + y + z = 6 & Loodvlak door $P$: $3 + 1 + 2 = 6$. \\
r + r + r = 6 \iff r = 2,\quad V(2, 2, 2) & Snijpunt met $l$. \\
|PV| = \sqrt{1 + 1 + 0} = \sqrt2 & De afstand.
\end{stappenbreed}}

\begin{oefeningen}
\oef{h4-afstand}

\oef{h4-afstandrechte}
\end{oefeningen}

%========================================================================================
\opdrachtenkop

\begin{oefeningen}
\needspace{14\baselineskip}
\oef{h4-o9}

\needspace{14\baselineskip}
\oef{h4-o11}

\oef{h4-o2}

\needspace{14\baselineskip}
\oef{h4-o12}

\needspace{12\baselineskip}
\oef{h4-o14}

\oef{h4-o1}

\needspace{12\baselineskip}
\oef{h4-o10}

\oef{h4-o3}

\needspace{16\baselineskip}
\oef{h4-o13}

\oef{h4-o15}

\oef{h4-o19}

\needspace{14\baselineskip}
\oef{h4-o17}

\oef{h4-o22}

\oef{h4-o4}

\oef{h4-o20}

\oef{h4-o18}

\oef{h4-o16}

\oef{h4-o5}

\oef{h4-o23}

\needspace{12\baselineskip}
\oef{h4-o21}

\oef{h4-o6}

\oef{h4-o26}

\needspace{14\baselineskip}
\oef{h4-o24}

\oef{h4-o7}

\needspace{14\baselineskip}
\oef{h4-o25}

\needspace{14\baselineskip}
\oef{h4-o27}

\oef{h4-o8}
\end{oefeningen}

%----------------------------------------------------------------------------------------
% KOPPEL: verbanden met andere leerstof
%----------------------------------------------------------------------------------------
\needspace{20\baselineskip}
\opdrachtsectie{Koppel}

\lkr{Goniometrische functies, Statistiek en Matrices: kijk na wat al gezien is. Het pakket legt de volgorde niet vast.}

\begin{oefeningen}
\oef{h4-koppel-1}

\oef{h4-koppel-2}

\oef{h4-koppel-3}
\end{oefeningen}

%----------------------------------------------------------------------------------------
\opdrachtsectie{Test jezelf}

\begin{oefeningen}
\oef{h4-t1}

\oef{h4-t6}

\oef{h4-t2}

\oef{h4-t5}

\oef{h4-t3}

\oef{h4-t7}

\oef{h4-t4}
\end{oefeningen}

\herhaalschema{%
  \fv{oef:h4-t1}{vb:h4-hoek}{oef:h4-hoek}%
  \fv{oef:h4-t6}{vb:h4-hoek}{oef:h4-scalair}%
  \fv{oef:h4-t2}{vb:h4-loodvlak}{oef:h4-loodvlak}%
  \fv{oef:h4-t5}{vb:h4-loodvlakken}{oef:h4-hoekvlakken}%
  \fv{oef:h4-t3}{vb:h4-afstandvlak}{oef:h4-afstandvlak}%
  \fv{oef:h4-t7}{vb:h4-afstandvlak}{oef:h4-o15}%
  \fv{oef:h4-t4}{vb:h4-afstandrechte}{oef:h4-afstandrechte}}

\ifzelfstudie
\section*{Maak je eigen samenvatting}
Sluit dit boek en schrijf in je schrift: de definitie van het scalair product en het verband met de hoek, de voorwaarde voor loodrechte stand, de normaalvector van een vlak, de vergelijking van een vlak door een punt met een gegeven normaalvector, de formule voor de afstand van een punt tot een vlak, en hoe je de afstand tot een rechte berekent. Vergelijk daarna met de kaders.
\fi
